\section{Results}

\subsection{First attempt}
At the first encounter with a question (\FirstN{} responses from \NEmployees{} employees), accuracy
was \FirstAcc\% (95\% CI \FirstAccCI; Table~\ref{tab:pierwsza_proba_wiek}). Accuracy varied with age
non-monotonically: it was lowest in the \FirstAccMinGroup{} age group (\FirstAccMin\%) and highest in
the \FirstAccMaxGroup{} age group (\FirstAccMax\%). When responses to the same questions were
compared, the odds of a correct response increased with age (OR \FirstAgeOR{} per 10 years, 95\% CI
\FirstAgeORCI; $p \FirstAgeP$). Without question fixed effects, the association was weaker
(OR \FirstAgeRawOR, 95\% CI \FirstAgeRawORCI; $p \FirstAgeRawP$).

\begin{table}[H]
\centering
\caption{Accuracy at the first attempt, overall and by employee age group.}
\label{tab:pierwsza_proba_wiek}
\footnotesize
\setlength{\tabcolsep}{4pt}
\tabelacsv{l*{7}{r}}{pierwsza_proba_wiek}%
\uwagatabeli{The first attempt is the first response in a sequence. 95\% CIs were obtained with the
bootstrap. The association with age was estimated by the logistic regression described in
Section~\ref{sec:analiza} (\FirstAgeN{} responses from \FirstAgeEmp{} employees of known age):
OR \FirstAgeOR{} per 10 years, 95\% CI \FirstAgeORCI, $p \FirstAgeP$.}
\end{table}

\subsection{Changes within a sequence (RQ1)}
Accuracy increased with attempt number, from \AccAttOne\% \AccAttOneCI{} at the first attempt to
\AccAttMax\% \AccAttMaxCI{} at attempt \AttMax{} and later (Figure~\ref{fig:krzywa_uczenia}). Because
the number of sequences decreased with each attempt, the change was examined in a cohort of
\CohortN{} sequences. In the cohort, accuracy rose from \CohortAccOne\% \CohortAccOneCI{} at
attempt~1 to \CohortAccK\% \CohortAccKCI{} at attempt~\CohortK{}, that is, by \SeqChange{} pp
(95\% CI \SeqChangeCI; Table~\ref{tab:sekwencje_wiek}). An increase within the same sequences means
that the improvement did not result from the earlier termination of sequences with lower accuracy.
Accuracy at attempt~1 in the cohort (\CohortAccOne\%) was close to accuracy at the first attempt in
all sequences (\FirstAcc\%).

The estimated increase was positive in all age groups and ranged from \SeqChangeMin{} pp
(\SeqChangeMinGroup{} years) to \SeqChangeMax{} pp (\SeqChangeMaxGroup{} years), although in the small
extreme age groups the lower limit of the 95\% CI was close to zero. The rate of increase did not
depend on age (interaction OR \SeqAgeOR{} per 10 years of age and doubling of attempt number, 95\% CI
\SeqAgeORCI; $p \SeqAgeP$).

\begin{figure}[H]
\centering
\includegraphics{\FigDir rys_krzywa_uczenia}
\caption{Accuracy by attempt number in all sequences available at a given attempt (their number
decreases with attempt number). Error bars: 95\% CI (bootstrap). Attempts \MaxAttemptPooled{} and
later were pooled. The course in the cohort is shown in Figure~\ref{fig:stany_kohorty}.}
\label{fig:krzywa_uczenia}
\end{figure}

\begin{table}[H]
\centering
\caption{Accuracy at successive attempts in the cohort, overall and by employee age group.}
\label{tab:sekwencje_wiek}
\footnotesize
\setlength{\tabcolsep}{4pt}
\tabelacsv{l*{7}{r}}{sekwencje_wiek}%
\uwagatabeli{Cohort: sequences of at least \CohortK{} attempts, attempts 1 to \CohortK{}.
pp: percentage points; 95\% CIs of the change were obtained with the bootstrap. The association between
the rate of improvement and age was estimated by the logistic regression described in
Section~\ref{sec:analiza} (\SeqAgeN{} responses from employees of known age).}
\end{table}

\subsection{Outcome of the previous attempt and interval (RQ2)}
The increase in accuracy results from two opposing processes: error correction and forgetting. The
schedule did not bias their comparison, because a question returned after a correct and after an
incorrect response with similar probability (\ReaskAfterCorrect\% and \ReaskAfterIncorrect\%) and
after a similar interval (median \GapAfterCorrect{} and \GapAfterIncorrect{} days). After a correct
response, the next one was correct in \RetainPct\% \RetainCI{} of cases, and forgetting occurred in
\ForgetPct\%. After an incorrect response, an error correction followed in \RecoverPct\%
\RecoverCI{} of cases. These rates refer to different groups of attempts, so the numbers of events were
also compared: there were \NCorrections{} error corrections and \NForgetting{} forgetting events. The
predominance of corrections explains the increase in accuracy. At constant error-correction and
forgetting rates, accuracy would tend toward an equilibrium level of \EquilibriumPct\%. Accuracy in the
cohort at attempt~\CohortK{} (\CohortAccK\%) was close to this level, but at attempts \AttMax{} and
later it was \AccAttMax\% and exceeded it. This is because the forgetting rate was not constant but
depended on the interval and on the number of consecutive correct responses.

The forgetting rate increased with the interval, from \ForgetShort\% at intervals shorter than 14 days
to \ForgetLong\% at intervals of at least 240 days (Figure~\ref{fig:odstep}). The error-correction
rate changed less with the interval (\RetShortIncorrect\% and \RetLongIncorrect\% in the same bins).

The forgetting rate also depended on the number of consecutive correct responses. After one correct
response, it was \ForgetRunOne\% \ForgetRunOneCI, after two \ForgetRunTwo\%, after three
\ForgetRunThree\%, and after at least \ForgetRunK{} \ForgetRunMax\% \ForgetRunMaxCI{}
(Table~\ref{tab:seria_poprawnych}). The difference persisted at every interval: at intervals shorter
than 14 days, the forgetting rate was \ForgetFreshShort\% after one and \ForgetStableShort\% after at
least three correct responses, and at intervals of at least 240 days \ForgetFreshLong\% and
\ForgetStableLong\%, respectively. Error corrections were the least durable. After a correction, an
error relapse occurred in \RelapsePct\% \RelapseCI{} of cases ($n = \RelapseN$), and in
\RelapseSamePct\% \RelapseSameCI{} of these the employee chose the same distractor as before the
correction. This share was similar to the percentage of repeated choices of the same distractor after
two successive errors (\SameDistractorPct\%; Section~\ref{sec:pb3}).

\begin{figure}[H]
\centering
\includegraphics{\FigDir rys_odstep}
\caption{Accuracy by interval, separately for attempts preceded by a correct response (circles) and by
an incorrect response (squares). For attempts preceded by a correct response, the forgetting rate is
100\% minus accuracy; for attempts preceded by an incorrect response, accuracy is the error-correction
rate. Error bars: 95\% CI (bootstrap).}
\label{fig:odstep}
\end{figure}

\begin{table}[H]
\centering
\caption{Forgetting rate by the number of consecutive correct responses given immediately before
a given attempt.}
\label{tab:seria_poprawnych}
\footnotesize
\setlength{\tabcolsep}{5pt}
\tabelacsv{l*{4}{r}}{seria_poprawnych}%
\uwagatabeli{Attempts preceded by a correct response. Forgetting rate: percentage of incorrect
responses. 95\% CIs were obtained with the bootstrap.}
\end{table}

\subsection{Persistent errors (RQ3)}
\label{sec:pb3}
Most errors were corrected: in \EverCorrectedPct\% of the \FirstWrongSeqN{} sequences that began with
an incorrect response and comprised at least two attempts, the employee later answered correctly at
least once. The chance of correction, however, decreased with each successive error. After one error,
accuracy was \CorrRunOne\% \CorrRunOneCI, after two consecutive errors \CorrRunTwo\%, after three
\CorrRunThree\%, and after at least \RunMax{} \CorrRunMax\% \CorrRunMaxCI{}
(Table~\ref{tab:seria_bledow}). Part of this decrease may be due to several consecutive errors
occurring mainly in more difficult questions.

Repeated errors were not consistent with random choice. After two successive errors, the employee
chose the same distractor in \SameDistractorPct\% \SameDistractorCI{} of cases
($n = \SameDistractorN$), whereas \SameDistractorExpected\% would be expected under random choice and
\SameDistractorExpectedPop\% under choice proportional to distractor popularity. This percentage did
not depend on the interval (\SameDistrGapShort\% at intervals shorter than 30 days,
\SameDistrGapMid\% at 30--89 days and \SameDistrGapLong\% at intervals of at least 90 days). Repeated errors were
concentrated in few questions: \PersistTopItems{} of the \PersistItems{} questions accounted for
\PersistTopItemsShare\% of them.

\begin{table}[H]
\centering
\caption{Accuracy by the number of consecutive errors made immediately before a given attempt.}
\label{tab:seria_bledow}
\footnotesize
\setlength{\tabcolsep}{5pt}
\tabelacsv{l*{5}{r}}{seria_bledow}%
\uwagatabeli{Attempts from the second onward. Column 0: attempts preceded by a correct response.
95\% CIs were obtained with the bootstrap.}
\end{table}

\subsection{Summary of results}
Figure~\ref{fig:stany_kohorty} summarizes the results in the cohort. Accuracy increases from attempt
to attempt (RQ1), error corrections outnumber forgetting events (RQ2), and at every attempt from the
second onward, part of the incorrect responses are repeated errors (RQ3).

\begin{figure}[H]
\centering
\includegraphics{\FigDir rys_stany_kohorty}
\caption{Accuracy and response state in the cohort (\CohortN{} sequences, attempts 1 to \CohortK{}).
Bars: share of responses by the outcome of the current and the previous attempt
(Section~\ref{sec:analiza}); at attempt~1, only correct and incorrect responses. Line: accuracy with
95\% CI (bootstrap), values above the bars.}
\label{fig:stany_kohorty}
\end{figure}
